% --- Archon physics formalization source begin ---
% archon:physics
% archon:covers PhyXMiniProblems/problem_phyx_mini_0456.lean
% archon:source-report reports/phyx_mini/problem_phyx_mini_0456.source.json

\chapter{Physics problem phyx\_mini\_0456}
\label{ch:PhyXMiniProblems_problem_phyx_mini_0456}

\paragraph{Problem source.}
\#\# Physical scenario

An engine, shown in figure, consists of a \$100 \textbackslash{}, \textbackslash{}text\{kg\}\$ cast iron block with a \$20 \textbackslash{}, \textbackslash{}text\{kg\}\$ aluminum head, \$20 \textbackslash{}, \textbackslash{}text\{kg\}\$ of steel parts, \$5 \textbackslash{}, \textbackslash{}text\{kg\}\$ of engine oil, and \$6 \textbackslash{}, \textbackslash{}text\{kg\}\$ of glycerine (antifreeze).

\#\# Question

All initial temperatures are \$5 \textbackslash{}, \textasciicircum{}\{\textbackslash{}circ\}\textbackslash{}text\{C\}\$, and as the engine starts we want to know how hot it becomes if it absorbs a net of \$7000 \textbackslash{}, \textbackslash{}text\{kJ\}\$ before it reaches a steady uniform temperature.

\#\# Answer choices

- A: A: 50 \textbackslash{}

- B: B: 60 \textbackslash{}

- C: C: 100 \textbackslash{}

- D: D: 80 \textbackslash{}

\#\# Figure caption (auxiliary; use the image as primary evidence)

The image illustrates a diagram of an automobile engine. It is a cross-sectional view that includes several key components:

1. **Cylinders and Pistons:** There are four pistons housed within cylinders, positioned in a V-shape, typical of a V-type engine. Each cylinder has a piston inside, connected to the crankshaft via connecting rods.

2. **Crankshaft:** This is located at the bottom of the engine, connected to the pistons and rotated as the pistons move.

3. **Camshaft and Valves:** At the top, there are components resembling a camshaft with rocker arms visible, controlling the opening and closing of the intake and exhaust valves.

4. **Timing Belt/Chain:** A V-shaped belt or chain is shown connecting the crankshaft and the camshaft, indicating the timing system.

5. **Cooling and Lubrication System:** Parts of cooling and lubrication systems, like a component that resembles an oil filter, are present.

6. **Label:** At the bottom, the text "Automobile engine" provides a label for the diagram.

The illustration is shaded in blue and gray tones, with blue lines highlighting certain mechanical components and pathways.

\#\# Dataset metadata

Specific Heat and Calorimetry; Physical Model Grounding Reasoning, Multi-Formula Reasoning

\paragraph{Recorded answer/context.}
D: D: 80 \textbackslash{}

\paragraph{Figure/image path.}
/root/proposal\_for\_physic/science-mango/phyx\_mini\_run/phyx\_data/test\_image/456.png

\paragraph{Formalization target.}
create a compiling Lean file with sorry bodies at `PhyXMiniProblems/problem\_phyx\_mini\_0456.lean`. The Lean declarations must preserve the physical quantities, dimensions or dimensional roles, figure labels, governing-law hypotheses, and final relation expressed by this problem.
Use Mathlib/Physlib names found through LeanExplore where available. If a physics API is missing, introduce faithful local abstractions rather than scalar placeholder aliases.

\begin{theorem}[Physics formalization target]
\label{thm:physics:phyx_mini_0456:target}
\lean{PhyXMiniProblems.ProblemPhyXMini0456.problem_phyx_mini_0456}
\uses{def:physics:phyx-mini-0456:phyxminiproblems-problemphyxmini0456-compositeengineheatingsetup, def:physics:phyx-mini-0456:phyxminiproblems-problemphyxmini0456-matchescompositeengineheatingscenario, def:physics:phyx-mini-0456:phyxminiproblems-problemphyxmini0456-matchesproblemreadouts, def:physics:phyx-mini-0456:phyxminiproblems-problemphyxmini0456-usesreferenceenginematerialthermaldata, def:physics:phyx-mini-0456:phyxminiproblems-problemphyxmini0456-matchessuppliedenginefigure, def:physics:phyx-mini-0456:phyxminiproblems-problemphyxmini0456-hasphysicalcompositeengineparameters, def:physics:phyx-mini-0456:phyxminiproblems-problemphyxmini0456-satisfiescompositeenginecalorimetrylaw, def:physics:phyx-mini-0456:phyxminiproblems-problemphyxmini0456-roundstodisplayedtemperature, def:physics:phyx-mini-0456:phyxminiproblems-problemphyxmini0456-isuniqueclosestdisplayedtemperature}
The assigned autoformalize agent should translate this physics problem into Lean declarations in the covered file, with theorem and lemma proof bodies written as `by sorry`.
\end{theorem}
\begin{proof}
This is an autoformalization task, not a proof task. Produce statements that can later be proved by the physics prover without weakening the physical model.
\end{proof}

% --- Archon declaration topology begin ---
\section{Declaration topology}
\begin{definition}[Mass Quantity]
  \label{def:physics:phyx-mini-0456:phyxminiproblems-problemphyxmini0456-massquantity}
  \lean{PhyXMiniProblems.ProblemPhyXMini0456.MassQuantity}
  A nonnegative physical mass, independent of the unit used to read it
\end{definition}
\begin{proof}
  This is the declared model definition of Mass Quantity.
\end{proof}
\begin{definition}[Specific Heat Capacity Quantity]
  \label{def:physics:phyx-mini-0456:phyxminiproblems-problemphyxmini0456-specificheatcapacityquantity}
  \lean{PhyXMiniProblems.ProblemPhyXMini0456.SpecificHeatCapacityQuantity}
  Mass-specific heat capacity, of dimension length² / (time² * temperature). Its SI readout is in joules per kilogram-kelvin
\end{definition}
\begin{proof}
  This is the declared model definition of Specific Heat Capacity Quantity.
\end{proof}
\begin{definition}[mass Readout]
  \label{def:physics:phyx-mini-0456:phyxminiproblems-problemphyxmini0456-massreadout}
  \lean{PhyXMiniProblems.ProblemPhyXMini0456.massReadout}
  \uses{def:physics:phyx-mini-0456:phyxminiproblems-problemphyxmini0456-massquantity}
  Read a physical mass in a selected mass unit
\end{definition}
\begin{proof}
  This is the declared model definition of mass Readout.
\end{proof}
\begin{definition}[mass In Kilograms]
  \label{def:physics:phyx-mini-0456:phyxminiproblems-problemphyxmini0456-massinkilograms}
  \lean{PhyXMiniProblems.ProblemPhyXMini0456.massInKilograms}
  \uses{def:physics:phyx-mini-0456:phyxminiproblems-problemphyxmini0456-massquantity, def:physics:phyx-mini-0456:phyxminiproblems-problemphyxmini0456-massreadout}
  Read a physical mass in kilograms
\end{definition}
\begin{proof}
  This is the declared model definition of mass In Kilograms.
\end{proof}
\begin{definition}[energy In Joules]
  \label{def:physics:phyx-mini-0456:phyxminiproblems-problemphyxmini0456-energyinjoules}
  \lean{PhyXMiniProblems.ProblemPhyXMini0456.energyInJoules}
  Read a physical energy in joules
\end{definition}
\begin{proof}
  This is the declared model definition of energy In Joules.
\end{proof}
\begin{definition}[energy In Kilojoules]
  \label{def:physics:phyx-mini-0456:phyxminiproblems-problemphyxmini0456-energyinkilojoules}
  \lean{PhyXMiniProblems.ProblemPhyXMini0456.energyInKilojoules}
  \uses{def:physics:phyx-mini-0456:phyxminiproblems-problemphyxmini0456-energyinjoules}
  Read a physical energy in kilojoules
\end{definition}
\begin{proof}
  This is the declared model definition of energy In Kilojoules.
\end{proof}
\begin{definition}[specific Heat In Joules Per Kilogram Kelvin]
  \label{def:physics:phyx-mini-0456:phyxminiproblems-problemphyxmini0456-specificheatinjoulesperkilogramkelvin}
  \lean{PhyXMiniProblems.ProblemPhyXMini0456.specificHeatInJoulesPerKilogramKelvin}
  \uses{def:physics:phyx-mini-0456:phyxminiproblems-problemphyxmini0456-specificheatcapacityquantity}
  Read a specific heat capacity in joules per kilogram-kelvin
\end{definition}
\begin{proof}
  This is the declared model definition of specific Heat In Joules Per Kilogram Kelvin.
\end{proof}
\begin{definition}[Celsius Temperature Reading]
  \label{def:physics:phyx-mini-0456:phyxminiproblems-problemphyxmini0456-celsiustemperaturereading}
  \lean{PhyXMiniProblems.ProblemPhyXMini0456.CelsiusTemperatureReading}
  Physlib's Temperature represents absolute temperature on a zero-preserving scale. This structure couples it to the affine Celsius readout printed in the exercise, with the absolute magnitude calibrated in kelvins
\end{definition}
\begin{proof}
  This is the declared model definition of Celsius Temperature Reading.
\end{proof}
\begin{definition}[Engine Component]
  \label{def:physics:phyx-mini-0456:phyxminiproblems-problemphyxmini0456-enginecomponent}
  \lean{PhyXMiniProblems.ProblemPhyXMini0456.EngineComponent}
  The five thermally significant engine components listed in the prose
\end{definition}
\begin{proof}
  This is the declared model definition of Engine Component.
\end{proof}
\begin{definition}[Thermal Material]
  \label{def:physics:phyx-mini-0456:phyxminiproblems-problemphyxmini0456-thermalmaterial}
  \lean{PhyXMiniProblems.ProblemPhyXMini0456.ThermalMaterial}
  The material assigned to each thermally significant component
\end{definition}
\begin{proof}
  This is the declared model definition of Thermal Material.
\end{proof}
\begin{definition}[material]
  \label{def:physics:phyx-mini-0456:phyxminiproblems-problemphyxmini0456-enginecomponent-material}
  \lean{PhyXMiniProblems.ProblemPhyXMini0456.EngineComponent.material}
  \uses{def:physics:phyx-mini-0456:phyxminiproblems-problemphyxmini0456-enginecomponent, def:physics:phyx-mini-0456:phyxminiproblems-problemphyxmini0456-thermalmaterial}
  The material composition specified for each engine component
\end{definition}
\begin{proof}
  This is the declared model definition of material.
\end{proof}
\begin{definition}[Figure Object]
  \label{def:physics:phyx-mini-0456:phyxminiproblems-problemphyxmini0456-figureobject}
  \lean{PhyXMiniProblems.ProblemPhyXMini0456.FigureObject}
  Mechanical objects visibly represented in the supplied cutaway raster
\end{definition}
\begin{proof}
  This is the declared model definition of Figure Object.
\end{proof}
\begin{definition}[Figure Label]
  \label{def:physics:phyx-mini-0456:phyxminiproblems-problemphyxmini0456-figurelabel}
  \lean{PhyXMiniProblems.ProblemPhyXMini0456.FigureLabel}
  Literal text labels visible in the supplied raster
\end{definition}
\begin{proof}
  This is the declared model definition of Figure Label.
\end{proof}
\begin{definition}[Cylinder Bank Arrangement]
  \label{def:physics:phyx-mini-0456:phyxminiproblems-problemphyxmini0456-cylinderbankarrangement}
  \lean{PhyXMiniProblems.ProblemPhyXMini0456.CylinderBankArrangement}
  Qualitative arrangement of the two cylinder banks in the raster
\end{definition}
\begin{proof}
  This is the declared model definition of Cylinder Bank Arrangement.
\end{proof}
\begin{definition}[Supplied Engine Figure]
  \label{def:physics:phyx-mini-0456:phyxminiproblems-problemphyxmini0456-suppliedenginefigure}
  \lean{PhyXMiniProblems.ProblemPhyXMini0456.SuppliedEngineFigure}
  \uses{def:physics:phyx-mini-0456:phyxminiproblems-problemphyxmini0456-figureobject, def:physics:phyx-mini-0456:phyxminiproblems-problemphyxmini0456-figurelabel, def:physics:phyx-mini-0456:phyxminiproblems-problemphyxmini0456-cylinderbankarrangement}
  Qualitative evidence retained from the primary image. The raster identifies the apparatus as an automobile engine and shows its cutaway geometry, but it contains no printed thermal datum or answer value
\end{definition}
\begin{proof}
  This is the declared model definition of Supplied Engine Figure.
\end{proof}
\begin{definition}[Composite Engine Heating Setup]
  \label{def:physics:phyx-mini-0456:phyxminiproblems-problemphyxmini0456-compositeengineheatingsetup}
  \lean{PhyXMiniProblems.ProblemPhyXMini0456.CompositeEngineHeatingSetup}
  \uses{def:physics:phyx-mini-0456:phyxminiproblems-problemphyxmini0456-massquantity, def:physics:phyx-mini-0456:phyxminiproblems-problemphyxmini0456-specificheatcapacityquantity, def:physics:phyx-mini-0456:phyxminiproblems-problemphyxmini0456-celsiustemperaturereading, def:physics:phyx-mini-0456:phyxminiproblems-problemphyxmini0456-enginecomponent, def:physics:phyx-mini-0456:phyxminiproblems-problemphyxmini0456-thermalmaterial, def:physics:phyx-mini-0456:phyxminiproblems-problemphyxmini0456-suppliedenginefigure}
  Independent physical observables and material properties of the heating experiment. The final uniform temperature is a free physical reading; it is not defined from the recorded answer or from 80 °C
\end{definition}
\begin{proof}
  This is the declared model definition of Composite Engine Heating Setup.
\end{proof}
\begin{definition}[final Temperature In Degrees Celsius]
  \label{def:physics:phyx-mini-0456:phyxminiproblems-problemphyxmini0456-finaltemperatureindegreescelsius}
  \lean{PhyXMiniProblems.ProblemPhyXMini0456.finalTemperatureInDegreesCelsius}
  \uses{def:physics:phyx-mini-0456:phyxminiproblems-problemphyxmini0456-compositeengineheatingsetup}
  Celsius readout of the common equilibrium temperature being requested
\end{definition}
\begin{proof}
  This is the declared model definition of final Temperature In Degrees Celsius.
\end{proof}
\begin{definition}[component Heat Capacity In Joules Per Kelvin]
  \label{def:physics:phyx-mini-0456:phyxminiproblems-problemphyxmini0456-componentheatcapacityinjoulesperkelvin}
  \lean{PhyXMiniProblems.ProblemPhyXMini0456.componentHeatCapacityInJoulesPerKelvin}
  \uses{def:physics:phyx-mini-0456:phyxminiproblems-problemphyxmini0456-massinkilograms, def:physics:phyx-mini-0456:phyxminiproblems-problemphyxmini0456-specificheatinjoulesperkilogramkelvin, def:physics:phyx-mini-0456:phyxminiproblems-problemphyxmini0456-enginecomponent, def:physics:phyx-mini-0456:phyxminiproblems-problemphyxmini0456-compositeengineheatingsetup}
  Heat capacity of one component in joules per kelvin
\end{definition}
\begin{proof}
  This is the declared model definition of component Heat Capacity In Joules Per Kelvin.
\end{proof}
\begin{definition}[total Heat Capacity In Joules Per Kelvin]
  \label{def:physics:phyx-mini-0456:phyxminiproblems-problemphyxmini0456-totalheatcapacityinjoulesperkelvin}
  \lean{PhyXMiniProblems.ProblemPhyXMini0456.totalHeatCapacityInJoulesPerKelvin}
  \uses{def:physics:phyx-mini-0456:phyxminiproblems-problemphyxmini0456-enginecomponent, def:physics:phyx-mini-0456:phyxminiproblems-problemphyxmini0456-compositeengineheatingsetup, def:physics:phyx-mini-0456:phyxminiproblems-problemphyxmini0456-componentheatcapacityinjoulesperkelvin}
  Sum of the five component heat capacities in joules per kelvin
\end{definition}
\begin{proof}
  This is the declared model definition of total Heat Capacity In Joules Per Kelvin.
\end{proof}
\begin{definition}[Matches Composite Engine Heating Scenario]
  \label{def:physics:phyx-mini-0456:phyxminiproblems-problemphyxmini0456-matchescompositeengineheatingscenario}
  \lean{PhyXMiniProblems.ProblemPhyXMini0456.MatchesCompositeEngineHeatingScenario}
  \uses{def:physics:phyx-mini-0456:phyxminiproblems-problemphyxmini0456-compositeengineheatingsetup}
  Qualitative process conditions explicitly described in the problem
\end{definition}
\begin{proof}
  This is the declared model definition of Matches Composite Engine Heating Scenario.
\end{proof}
\begin{definition}[Matches Problem Readouts]
  \label{def:physics:phyx-mini-0456:phyxminiproblems-problemphyxmini0456-matchesproblemreadouts}
  \lean{PhyXMiniProblems.ProblemPhyXMini0456.MatchesProblemReadouts}
  \uses{def:physics:phyx-mini-0456:phyxminiproblems-problemphyxmini0456-massinkilograms, def:physics:phyx-mini-0456:phyxminiproblems-problemphyxmini0456-energyinkilojoules, def:physics:phyx-mini-0456:phyxminiproblems-problemphyxmini0456-compositeengineheatingsetup}
  The five masses, common initial temperature, and absorbed energy printed in the problem. No field mentions the final temperature or an answer choice
\end{definition}
\begin{proof}
  This is the declared model definition of Matches Problem Readouts.
\end{proof}
\begin{definition}[Uses Reference Engine Material Thermal Data]
  \label{def:physics:phyx-mini-0456:phyxminiproblems-problemphyxmini0456-usesreferenceenginematerialthermaldata}
  \lean{PhyXMiniProblems.ProblemPhyXMini0456.UsesReferenceEngineMaterialThermalData}
  \uses{def:physics:phyx-mini-0456:phyxminiproblems-problemphyxmini0456-specificheatinjoulesperkilogramkelvin, def:physics:phyx-mini-0456:phyxminiproblems-problemphyxmini0456-compositeengineheatingsetup}
  Standard reference values for the five material specific heats. These are supplemental textbook data required by the numerical question rather than values printed in the problem. None is a final-temperature assertion
\end{definition}
\begin{proof}
  This is the declared model definition of Uses Reference Engine Material Thermal Data.
\end{proof}
\begin{definition}[Matches Supplied Engine Figure]
  \label{def:physics:phyx-mini-0456:phyxminiproblems-problemphyxmini0456-matchessuppliedenginefigure}
  \lean{PhyXMiniProblems.ProblemPhyXMini0456.MatchesSuppliedEngineFigure}
  \uses{def:physics:phyx-mini-0456:phyxminiproblems-problemphyxmini0456-suppliedenginefigure}
  Qualitative facts read from the primary automobile-engine raster
\end{definition}
\begin{proof}
  This is the declared model definition of Matches Supplied Engine Figure.
\end{proof}
\begin{definition}[Has Physical Composite Engine Parameters]
  \label{def:physics:phyx-mini-0456:phyxminiproblems-problemphyxmini0456-hasphysicalcompositeengineparameters}
  \lean{PhyXMiniProblems.ProblemPhyXMini0456.HasPhysicalCompositeEngineParameters}
  \uses{def:physics:phyx-mini-0456:phyxminiproblems-problemphyxmini0456-massinkilograms, def:physics:phyx-mini-0456:phyxminiproblems-problemphyxmini0456-energyinjoules, def:physics:phyx-mini-0456:phyxminiproblems-problemphyxmini0456-specificheatinjoulesperkilogramkelvin, def:physics:phyx-mini-0456:phyxminiproblems-problemphyxmini0456-compositeengineheatingsetup}
  Positivity conditions selecting a physical heating experiment
\end{definition}
\begin{proof}
  This is the declared model definition of Has Physical Composite Engine Parameters.
\end{proof}
\begin{definition}[Satisfies Composite Engine Calorimetry Law]
  \label{def:physics:phyx-mini-0456:phyxminiproblems-problemphyxmini0456-satisfiescompositeenginecalorimetrylaw}
  \lean{PhyXMiniProblems.ProblemPhyXMini0456.SatisfiesCompositeEngineCalorimetryLaw}
  \uses{def:physics:phyx-mini-0456:phyxminiproblems-problemphyxmini0456-energyinjoules, def:physics:phyx-mini-0456:phyxminiproblems-problemphyxmini0456-enginecomponent, def:physics:phyx-mini-0456:phyxminiproblems-problemphyxmini0456-compositeengineheatingsetup, def:physics:phyx-mini-0456:phyxminiproblems-problemphyxmini0456-finaltemperatureindegreescelsius, def:physics:phyx-mini-0456:phyxminiproblems-problemphyxmini0456-componentheatcapacityinjoulesperkelvin}
  The governing composite-body calorimetry law Q = ∑ᵢ mᵢ cᵢ (T\_final - T\_initial,i). Because Celsius and kelvin increments have the same size, the calibrated Celsius differences may be used with SI specific heats. The equation relates the unknown final temperature to the input data but does not assign it the requested numerical value
\end{definition}
\begin{proof}
  This is the declared model definition of Satisfies Composite Engine Calorimetry Law.
\end{proof}
\begin{definition}[Answer Choice]
  \label{def:physics:phyx-mini-0456:phyxminiproblems-problemphyxmini0456-answerchoice}
  \lean{PhyXMiniProblems.ProblemPhyXMini0456.AnswerChoice}
  Labels of the four temperatures printed in the answer list
\end{definition}
\begin{proof}
  This is the declared model definition of Answer Choice.
\end{proof}
\begin{definition}[displayed Temperature In Degrees Celsius]
  \label{def:physics:phyx-mini-0456:phyxminiproblems-problemphyxmini0456-displayedtemperatureindegreescelsius}
  \lean{PhyXMiniProblems.ProblemPhyXMini0456.displayedTemperatureInDegreesCelsius}
  \uses{def:physics:phyx-mini-0456:phyxminiproblems-problemphyxmini0456-answerchoice}
  Celsius values printed beside the four answer labels
\end{definition}
\begin{proof}
  This is the declared model definition of displayed Temperature In Degrees Celsius.
\end{proof}
\begin{definition}[recorded Dataset Answer]
  \label{def:physics:phyx-mini-0456:phyxminiproblems-problemphyxmini0456-recordeddatasetanswer}
  \lean{PhyXMiniProblems.ProblemPhyXMini0456.recordedDatasetAnswer}
  \uses{def:physics:phyx-mini-0456:phyxminiproblems-problemphyxmini0456-answerchoice}
  The answer label recorded by the source dataset, retained as metadata
\end{definition}
\begin{proof}
  This is the declared model definition of recorded Dataset Answer.
\end{proof}
\begin{definition}[Rounds To Displayed Temperature]
  \label{def:physics:phyx-mini-0456:phyxminiproblems-problemphyxmini0456-roundstodisplayedtemperature}
  \lean{PhyXMiniProblems.ProblemPhyXMini0456.RoundsToDisplayedTemperature}
  \uses{def:physics:phyx-mini-0456:phyxminiproblems-problemphyxmini0456-compositeengineheatingsetup, def:physics:phyx-mini-0456:phyxminiproblems-problemphyxmini0456-finaltemperatureindegreescelsius, def:physics:phyx-mini-0456:phyxminiproblems-problemphyxmini0456-answerchoice, def:physics:phyx-mini-0456:phyxminiproblems-problemphyxmini0456-displayedtemperatureindegreescelsius}
  The final reading rounds to a displayed whole-ten Celsius value
\end{definition}
\begin{proof}
  This is the declared model definition of Rounds To Displayed Temperature.
\end{proof}
\begin{definition}[Is Closest Displayed Temperature]
  \label{def:physics:phyx-mini-0456:phyxminiproblems-problemphyxmini0456-isclosestdisplayedtemperature}
  \lean{PhyXMiniProblems.ProblemPhyXMini0456.IsClosestDisplayedTemperature}
  \uses{def:physics:phyx-mini-0456:phyxminiproblems-problemphyxmini0456-compositeengineheatingsetup, def:physics:phyx-mini-0456:phyxminiproblems-problemphyxmini0456-finaltemperatureindegreescelsius, def:physics:phyx-mini-0456:phyxminiproblems-problemphyxmini0456-answerchoice, def:physics:phyx-mini-0456:phyxminiproblems-problemphyxmini0456-displayedtemperatureindegreescelsius}
  A displayed choice is at least as close as every displayed temperature
\end{definition}
\begin{proof}
  This is the declared model definition of Is Closest Displayed Temperature.
\end{proof}
\begin{definition}[Is Unique Closest Displayed Temperature]
  \label{def:physics:phyx-mini-0456:phyxminiproblems-problemphyxmini0456-isuniqueclosestdisplayedtemperature}
  \lean{PhyXMiniProblems.ProblemPhyXMini0456.IsUniqueClosestDisplayedTemperature}
  \uses{def:physics:phyx-mini-0456:phyxminiproblems-problemphyxmini0456-compositeengineheatingsetup, def:physics:phyx-mini-0456:phyxminiproblems-problemphyxmini0456-answerchoice, def:physics:phyx-mini-0456:phyxminiproblems-problemphyxmini0456-isclosestdisplayedtemperature}
  The selected choice is the unique closest displayed temperature
\end{definition}
\begin{proof}
  This is the declared model definition of Is Unique Closest Displayed Temperature.
\end{proof}
% --- Archon declaration topology end ---
% --- Archon physics formalization source end ---
